\section{The CPU reference walker}
\label{sec:cpuwalk}

The CPU walker is \code{OctreeTriangleTessellator::iterateTriangles}.  It is
no longer used at runtime, but it defines the semantics the GPU shader must
reproduce and it is the only mesher used by the offline Blender exporter.
Reading it is the best way to understand the GPU port because the shader is
a line-by-line translation with an explicit stack instead of C++
recursion.

\subsection{Frontier and scan cells}

The walk starts at a chunk node and a target LoD $r$.  It descends to the
cells at that rung:
\begin{lstlisting}[caption={Frontier descent (conceptual).},label={lst:frontier}]
void walk(node, cube, targetLod) {
    if (node.lod == targetLod) {
        scanCell(node, cube);
    } else if (node.lod > targetLod) {
        for (child in node.children) walk(child, cube.child(i), targetLod);
    }
}
\end{lstlisting}
A \emph{scan cell} is a cube at the target rung.  Only cells that straddle
the surface need to be scanned.  The GPU flattening compacts the list with
the same ``has negative and non-negative corner'' test, because a cell with
no sign change on any of its twelve edges cannot emit a triangle.

\subsection{Twelve edges and edge spans}

For a scan cell with cube minimum $\mathbf{c}$ and size $L$, each of the
twelve edges is converted to an \code{EdgeSpan}:
\begin{itemize}
\item \code{axis}: the edge's varying axis (the one with the largest
  coordinate difference, which is exact for an axis-aligned cube);
\item \code{fixedU}, \code{fixedV}: the two constant coordinates;
\item \code{start}, \code{end}: the varying coordinate interval;
\item \code{eps}: the geometric epsilon
  $\max(\text{worldSize}, L)\cdot 10^{-6}$;
\item \code{sampleOffset}: $L\cdot10^{-3}$, used to nudge boundary samples.
\end{itemize}
The span is active only when its two corner SDF values have opposite signs.
The sign of the endpoint is also remembered; it is the authoritative
winding oracle.

\subsection{Why one edge can produce many segments}

In a uniform grid, an active edge is shared by exactly four cells, all of
the same size.  The dual polygon is a quad (or a triangle when one of the
four cells is not a surface cell).  In an adaptive octree, the four
quadrant cells around an edge can have different sizes and can be nested.
The edge can cross several cell boundaries along its length.  The correct
mesh must split the edge at every such boundary; otherwise one side of the
boundary would connect to a vertex that the other side does not know.

\begin{figure}[htbp]
\centering
\begin{tikzpicture}[scale=0.9]
  % coarse cell left, fine right
  \draw[grid] (0,0) rectangle (3,3);
  \draw[grid] (3,0) rectangle (6,3);
  \draw[grid] (3,1.5) -- (6,1.5);
  \draw[grid] (4.5,0) -- (4.5,1.5);
  \draw[grid] (4.5,1.5) -- (4.5,3);
  % edge
  \draw[accent2,line width=1.2pt] (0,2.2) -- (6,2.2);
  % cell boundary crossings
  \foreach \x/\lab in {3.0/coarse|fine, 4.5/fine|fine} {
    \draw[accent,fill=white] (\x,2.2) circle (2pt);
  }
  \node[font=\scriptsize,above] at (3,2.25) {cell boundary};
  \node[font=\scriptsize,above] at (4.5,2.25) {cell boundary};
  % subsegments
  \draw[decorate,decoration={brace,amplitude=4pt},good]
    (0,1.85) -- (3,1.85) node[midway,below=4pt,font=\scriptsize]{segment A};
  \draw[decorate,decoration={brace,amplitude=4pt},good]
    (3,1.85) -- (4.5,1.85) node[midway,below=4pt,font=\scriptsize]{B};
  \draw[decorate,decoration={brace,amplitude=4pt},good]
    (4.5,1.85) -- (6,1.85) node[midway,below=4pt,font=\scriptsize]{C};
  % cell labels
  \node[font=\scriptsize] at (1.5,0.4) {coarse cell};
  \node[font=\scriptsize] at (3.75,0.4) {fine};
  \node[font=\scriptsize] at (5.25,2.7) {fine};
\end{tikzpicture}
\caption{An active edge crossing a coarse/fine cell boundary.  The edge is
split at every cell boundary coordinate; each sub-segment is emitted by its
owner cell.  Both the coarse and the fine side collect the same set of
boundaries, so the resulting rings agree along the shared edge.}
\label{fig:edge-split}
\end{figure}

\subsection{Collecting breaks}

The break collector is the heart of the adaptive stitching.  Given an edge
span and an interval $[a,b]$, it samples the four quadrant cells at the
midpoint.  Each valid cell contributes its cube's minimum and maximum
coordinate along the edge axis.  If these coordinates lie strictly inside
$[a,b]$, they become breakpoints.  The interval is then recursively split
at each breakpoint and the process repeats.  The recursion stops when the
interval is shorter than $2\epsilon$ or contains no new break.

The crucial detail is that the four quadrant cells are resolved by
\code{findCellAt}, a root-to-cell descent that mirrors the octree's
\code{getChildIndex} arithmetic.  A sample is nudged by one ulp in the two
fixed axes so that a point exactly on a cell face does not land in both
cells; a small normal offset replaces the ulp when the value is zero (the
next float is subnormal and may be flushed to zero).  The four samples are:
\begin{align}
  (\text{fixedU}-\delta,\;& \text{fixedV}-\delta),\\
  (\text{fixedU}+\delta,\;& \text{fixedV}-\delta),\\
  (\text{fixedU}-\delta,\;& \text{fixedV}+\delta),\\
  (\text{fixedU}+\delta,\;& \text{fixedV}+\delta).
\end{align}
The first sample is resolved with a full descent; the other three reuse it
if they still lie inside its cube, which is exact because \code{findCellAt}
is a pure function of position.  This is the same optimization in the CPU
walker and the shader.

\subsection{The break list and the interval stack}

The CPU collector uses recursive calls and a \code{std::vector<float>}
break list.  The GPU port replaces recursion with an explicit LIFO stack:
\begin{lstlisting}[style=glsl,caption={The GPU break collector (shape).},label={lst:breaks}]
void collectBreaks(edge, start, end) {
    pushTask(start, end);
    while (top > 0) {
        pop(a, b);
        if (b - a <= eps*2) continue;
        mid = 0.5*(a+b);
        local = {a,b};
        cells = sampleQuadrants(edge, mid);
        for each valid cell:
            addBreakLocal(cell.min[axis], a, b);
            addBreakLocal(cell.min[axis]+cell.size, a, b);
        sortUniqueLocal();
        for interior local breaks: addBreakGlobal(...);
        push sub-intervals in reverse order;
    }
}
\end{lstlisting}
The CPU version appends all interior breaks before recursing into the
sub-intervals.  The GPU version reproduces that order by pushing the
sub-intervals in reverse so the LIFO pop visits them left to right.  The
global break list is therefore identical up to floating-point ordering.

The fixed capacities are part of the contract:
\begin{center}
\begin{tabular}{@{}lll@{}}
\toprule
Buffer & Capacity & Meaning\\
\midrule
\code{BREAK\_CAP} & 64 & final breaks per edge span\\
\code{LOCAL\_BREAK\_CAP} & 32 & breaks inside one interval\\
\code{TASK\_STACK} & 256 & pending intervals\\
\code{POS\_MEMO} & 8 & QEF position cache per invocation\\
\bottomrule
\end{tabular}
\end{center}
A cap overflow does not truncate silently.  \code{noteOverflow()} sets a
per-invocation flag and increments \code{overflowCount}; the host fails the
chunk and retries it instead of drawing a hole.

\subsection{Owner selection}

For each sub-segment $[a,b]$, the owner is the surface quadrant with the
smallest ordering key:
\begin{enumerate}
\item smaller cube size first;
\item then smaller $x$ minimum;
\item then smaller $y$ minimum;
\item then smaller $z$ minimum;
\item then smaller node index (a deterministic stand-in for pointer order).
\end{enumerate}
The owner is computed from the four quadrants at the segment midpoint, not
from the scan cell.  This is what makes two neighbouring chunks agree: both
see the same four quadrant cells and both choose the same owner.

\subsection{Polygon reconstruction}

The four quadrants are filtered to surface cells and deduplicated by node
identity and by position (within the edge epsilon).  A coarse ancestor can
appear in two opposite quadrants of the ring; the old code only removed
adjacent duplicates and then rejected the whole segment, which produced
the missing surface along the world axes.  The current code deduplicates
against every already-collected quadrant.  If the resulting list has three
or four entries, it is rotated so the owner is first; then it is fanned into
one or two triangles.

\begin{figure}[htbp]
\centering
\begin{tikzpicture}[scale=0.95]
  \begin{scope}[shift={(0,0)}]
    \draw[grid] (0,0) rectangle (3,3);
    \draw[grid] (1,0)--(1,3); \draw[grid] (2,0)--(2,3);
    \draw[grid] (0,1)--(3,1); \draw[grid] (0,2)--(3,2);
    \fill[accent!18] (1,1)--(2,1)--(2,2)--(1,2)--cycle;
    \node[font=\scriptsize] at (1.5,1.5) {owner};
    \node[font=\scriptsize] at (0.5,1.5) {A};
    \node[font=\scriptsize] at (2.5,1.5) {B};
    \node[font=\scriptsize] at (1.5,2.5) {C};
    \node[font=\scriptsize] at (1.5,0.5) {D};
    \draw[accent2,line width=1pt] (0.5,1.5)--(1.5,1.5)--(2.5,1.5);
    \node[below,font=\scriptsize] at (1.5,-0.2) {edge midpoints};
  \end{scope}
  \begin{scope}[shift={(4.6,0)}]
    \draw[grid] (0,0) rectangle (3,3);
    \draw[accent2,fill=accent2!15] (1.0,0.8)--(2.2,1.0)--(2.0,2.2)--(0.8,2.0)--cycle;
    \draw[accent2,thick] (1.0,0.8)--(2.2,1.0)--(2.0,2.2)--(0.8,2.0)--cycle;
    \node[font=\scriptsize] at (1.5,1.5) {ring};
    \draw[accent2,thick] (0.8,2.0)--(2.0,2.2);
    \draw[accent2,thick] (2.0,2.2)--(2.2,1.0);
    \draw[accent2,thick] (2.2,1.0)--(1.0,0.8);
    \node[below,font=\scriptsize] at (1.5,-0.2) {winding from SDF sign};
  \end{scope}
\end{tikzpicture}
\caption{Left: the four quadrant cells around an edge.  Right: the dual
polygon is the ring of their vertices.  The ring is rotated so the owner
leads, then fanned.  The winding is chosen from the SDF sign at the edge
start, not from the normals.}
\label{fig:ring}
\end{figure}

\subsection{Winding}

The triangle order must be consistent with the SDF sign.  The walker takes
the interpolated SDF at the segment start using the owner's field.  If the
owner's interpolation is infinite or degenerates, it falls back to the
producing cell's endpoint sign.  A solid start means the ring is ordered
so the outward normal points away from the solid; a positive start means
the order is reversed.  The sign is an \emph{authoritative oracle}: it is
independent of the QEF normal, which can degrade at coarse levels and on
non-planar quads.  This is why the surface does not flip at LoD
transitions.

\subsection{Triangle fan and duplicate rejection}

A 3-ring produces one triangle; a 4-ring produces two, fanned around the
owner.  Before emission, the three dual positions are compared; if any two
are within the edge epsilon, the triangle is dropped.  A degenerate QEF
vertex (or two cells that collapsed to the same position) must not create a
zero-area triangle: it would break normals, shadows, and ray tracing.

\subsection{The CPU triangle handler}

The CPU walker passes each triangle to \code{Tesselator::handle}, which:
\begin{itemize}
\item resolves the three cell vertices to \code{Vertex} records;
\item copies brush, HSV, texture type, and mapping data;
\item computes the triplanar UV plane from the face normal;
\item appends the triangle to the chunk's \code{Geometry}.
\end{itemize}
The GPU port does not have a \code{Tesselator} equivalent.  Stage 1 writes
the vertex payload directly from the node id and the index stream directly
from the raw keys.  This is possible because the CPU handler's output is a
pure function of the node state.

\subsection{Why keep the CPU walker?}

The CPU walker is retained for three reasons:
\begin{enumerate}
\item \textbf{Offline export.}  The Blender exporter has no GPU and must
  produce the same mesh.
\item \textbf{Reference.}  The GPU shader's arithmetic can be compared
  chunk-by-chunk with the CPU walker under a self-check flag.
\item \textbf{Debugging.}  When a GPU mesh is wrong, the CPU walker can
  generate the expected triangles for a small chunk and the two sets can be
  diffed.
\end{enumerate}
It is not a fallback: the runtime never invokes it.  A chunk that cannot be
flattened is skipped, not CPU-meshed.
